% from ega2-1-fr.tex
\oldpage[II]{70}

\begin{proof}
Compte tenu de (\hyperref[II.3.8.3-fr]{3.8.3}), il suffit de démontrer que si
$r_{\mathcal{L},\psi}$ est partout définie et est une immersion, il en est de
même de $r_{\mathcal{L},\psi_\lambda}$ pour un $\lambda$ au moins. Par le
même raisonnement que dans (\hyperref[II.3.7.7-fr]{3.7.7}), utilisant la
quasi-compacité de $Y$, on se ramène d'abord au cas où $Y$ est affine. Comme
$X$ est quasi-compact, (\hyperref[II.3.8.2-fr]{3.8.2}) montre l'existence
d'une famille \emph{finie} $(s_i)$ d'éléments de $S$ ($s_i\in S_{n_i}$)
ayant les propriétés (i), (ii) et (iii). Le morphisme $X_{f_i}\to Y$ (où
$f_i=\psi^\flat(s_i)$) est de type fini~: en effet, c'est un morphisme de
schémas affines, donc il est quasi-compact
(\hyperref[I.6.6.1-fr]{I, 6.6.1}), et localement de type fini puisque $q$ est
de type fini (\hyperref[I.6.3.2-fr]{I, 6.3.2}), et la conclusion résulte de
(\hyperref[I.6.6.3-fr]{I, 6.6.3}). L'anneau $B_i$ de $X_{f_i}$ est par suite
une $A$-algèbre de type fini (\hyperref[I.6.3.3-fr]{I, 6.3.3})~; soit
$(t_{ij})$ une famille de générateurs de cette algèbre. Il y a par hypothèse
des éléments $s'_{ij}\in S_{m_{ij}n_i}$ tels que
\[
  t_{ij}=(\psi^\flat(s'_{ij})|X_{f_i})
  (\psi^\flat(s_i)|X_{f_i})^{-m_{ij}}.
\]
Il y a par hypothèse un $\lambda$ et des éléments
$s_{i\lambda}\in S_{n_i}^\lambda$,
$s'_{ij\lambda}\in S_{m_{ij}n_i}^\lambda$ dont les images par
$\varphi_\lambda$ sont respectivement $s_i$ et $s'_{ij}$~; il est clair que
$r_{\mathcal{L},\psi_\lambda}$ vérifie les conditions (i), (ii) et (iii) de
(\hyperref[II.3.8.2-fr]{3.8.2}).
\end{proof}

\begin{proposition}[3.8.5]
\label{II.3.8.5-fr}
Soient $Y$ un schéma quasi-compact, ou un préschéma dont l'espace sous-jacent
est noethérien, $q:X\to Y$ un morphisme de type fini, $\mathcal{L}$ un
$\mathcal{O}_X$-Module inversible, $\mathcal{S}$ une $\mathcal{O}_Y$-Algèbre
graduée quasi-cohérente,
$\psi:q^*(\mathcal{S})\to\bigoplus_{n\geq0}\mathcal{L}^{\otimes n}$ un
homomorphisme d'Algèbres graduées. Pour que $r_{\mathcal{L},\psi}$ soit
partout défini et soit une immersion, il faut et il suffit qu'il existe un
entier $n>0$ et un sous-$\mathcal{O}_Y$-Module quasi-cohérent de type fini
$\mathcal{E}$ de $\mathcal{S}_n$, tels que~:
\begin{enumerate}
  \item[a)] l'homomorphisme
    $\psi_n\circ q^*(j_n):q^*(\mathcal{E})\to\mathcal{L}^{\otimes n}$
    (où $j_n:\mathcal{E}\to\mathcal{S}$ est l'injection canonique) soit
    surjectif~;
  \item[b)] si on désigne par $\mathcal{S}'$ la sous-$\mathcal{O}_Y$-Algèbre
    (graduée) de $\mathcal{S}$ engendrée par $\mathcal{E}$, par $\psi'$
    l'homomorphisme $\psi\circ q^*(j')$ ($j'$ injection de $\mathcal{S}'$
    dans $\mathcal{S}$), $r_{\mathcal{L},\psi'}$ soit partout défini et soit
    une immersion.
\end{enumerate}

Lorsqu'il en est ainsi, tout sous-$\mathcal{O}_Y$-Module quasi-cohérent
$\mathcal{E}'$ de $\mathcal{S}_n$, contenant $\mathcal{E}$, possède la même
propriété, et il en est de même de l'image de $\bigotimes^k\mathcal{E}$ dans
$\mathcal{S}_{kn}$ pour tout $k>0$.
\end{proposition}

\begin{proof}
La suffisance de la condition et les deux dernières assertions sont des cas
particuliers de (\hyperref[II.3.8.3-fr]{3.8.3}), compte tenu de l'isomorphie
canonique entre $\operatorname{Proj}(\mathcal{S})$ et
$\operatorname{Proj}(\mathcal{S}^{(k)})$
(\hyperref[II.3.1.8-fr]{3.1.8}).

Soit $(U_i)$ un recouvrement ouvert fini de $Y$ par des ouverts affines et
posons $A_i=A(U_i)$. Comme $q^{-1}(U_i)$ est compact, l'hypothèse que
$r_{\mathcal{L},\psi}$ est une immersion partout définie et
(\hyperref[II.3.8.2-fr]{3.8.2}) entraînent l'existence d'une famille finie
$(s_{ij})$ d'éléments de
$S^{(i)}=\Gamma(U_i,\mathcal{S})$ ($s_{ij}\in S_{n_{ij}}^{(i)}$) ayant les
propriétés (i), (ii), (iii). On voit comme dans la démonstration de
(\hyperref[II.3.8.4-fr]{3.8.4}) que le morphisme $X_{f_{ij}}\to U_i$ (où
$f_{ij}=\psi^\flat(s_{ij})$) est de type fini, et l'anneau $B_{ij}$ de
$X_{f_{ij}}$ est par suite une $A_i$-algèbre de type fini, ayant un système de
générateurs de la forme
$(\psi^\flat(t_{ijk})|X_{f_{ij}})(f_{ij}|X_{f_{ij}})^{-m_{ijk}}$, où
$t_{ijk}\in S_{m_{ijk}n_{ij}}^{(i)}$. Soit $n$ un multiple commun des
$m_{ijk}n_{ij}$~; remplaçant (pour chaque $(i,j,k)$) $s_{ij}$ par une
puissance $s_{ij}^\rho$ telle que $\rho m_{ijk}n_{ij}=n$, et multipliant
$t_{ijk}$ par $s_{ij}^{\rho-m_{ijk}}$, on peut supposer que pour chaque $i$
tous les $s_{ij}$ et $t_{ijk}$ appartiennent à $S_n^{(i)}$ et que
$m_{ijk}=1$. Soit $E_i$ le sous-$A_i$-Module de $S_n^{(i)}$ engendré par ces
éléments (pour $i$ fixé). Il existe un sous-$\mathcal{O}_Y$-Module cohérent de
type

